\section{The complete small symplectic core reduction}
\label{sec:sp43-core}

\paragraph{Status.}
This new complete author derivation and its exact certificate have
passed the root reviewer's whole-proof and complete-code review.
The root reviewer also freshly reran the standalone certificate:
all $59$ selected rows and $1992$ nonscalar spectra passed.
That rerun is distinguished from a separate arithmetic implementation.
The imported Fourier routine had already passed the full prior
SPOR and A1 code reviews.  This section concerns every
irreducible representation of the indicated core, not just its Weil
constituents.  Its transitive-factor hypothesis is precisely the
one supplied by D-BI-19; no statement about arbitrary mixed
tensor products is inferred.

\subsection*{D-SPW-04a. Complete scalar-normalizer spectra}

Put $S=\operatorname{PSp}_4(3)\simeq\operatorname{PSU}_4(2)$,
and let $U=\operatorname{Sp}_4(3)$ be its universal perfect cover.
The multiplier and full outer group of $S$ both have order two;
see the published-cover input in D-SPW-02 and the official
CTblLib extension metadata, interpreted as in D-SPOR-01.
The relevant complete character tables have identifiers
\texttt{2.U4(2)} and \texttt{2.U4(2).2}.  Their orders are
$51840$ and $103680$ and their class counts are $34$ and $38$.
Both retained source strings identify the published ATLAS input.

\begin{lemma}[D-SPW-04a: every compatible normalizer]
\label{D-SPW-04a}
Let $\rho$ be any nonlinear complex irreducible representation
of $U$, of degree $a$, and let $T$ be a nonscalar finite-order
linear normalizer of its image.  Then
\[
 \nu(T)\ge
 \begin{cases}
 a/4,&a=4\text{ or }a\ge7,\\
 a/5,&a=5,\\
 a/6,&a=6.
 \end{cases}                                               \tag{SPW.12}
\]
There are no nonlinear degrees below four.  There is exactly
one degree-six core character; it has a trivial central
character and Frobenius--Schur indicator $+1$.

The degree-five quarter-bound expectation is false: both
degree-five characters have, on core class \texttt{2b}, the
involution spectrum $1^1,(-1)^4$, so $\nu=1$.
\end{lemma}

\begin{proof}
The compact exact input retains every character, every class
size and order, every prime power map, all source strings, and
the actual fusion of $U$ into its specified outer extension
$A=\texttt{2.U4(2).2}$.  The GAP exporter checks full row
orthogonality and the sum-of-squared-degrees identity.  These
are consistency checks; the group and character-table
identifications remain the stated published mathematical inputs.

For an element $x$ of order $d$, and each $0\le b<d$, the
exact eigenvalue multiplicity is
\[
 m_b=\frac1d\sum_{j=0}^{d-1}
               \chi(x^j)\zeta_d^{-bj}.                    \tag{SPW.13}
\]
The prime power maps and Galois conjugation determine every
$\chi(x^j)$.  The standalone checker evaluates (SPW.13) by
integer-polynomial reduction modulo $\Phi_d$, checks that
every result is a nonnegative integer, and checks
$\sum_bm_b=a$.  It selects the extension characters again by
pointwise equality of their restriction with a full irreducible
core-character row; thus the selection does not assume
irreducibility merely from degree.

There are $33$ nonlinear core rows and $26$ extension rows
with irreducible nonlinear restriction.  All $1992$ nonscalar
entries of those $59$ rows are checked.  The only entries
violating $\nu\ge a/4$ are
\[
\begin{array}{c|c|c|c|c}
\text{table}&\text{row}&a&\text{class}&\text{multiplicities}\\\hline
\texttt{2.U4(2)}&2,3&5&\texttt{2b},\ d=2&(1,4)\\
\texttt{2.U4(2).2}&4&6&\texttt{4d},\ d=4&(5,0,1,0)\\
\texttt{2.U4(2).2}&5&6&\texttt{4d},\ d=4&(1,0,5,0).
\end{array}                                                \tag{SPW.14}
\]
All multiplicities in (SPW.14) are ordered starting with the
eigenvalue one.  The two extension rows restrict to core row
four, the unique degree-six character.  That core character
has integer values, is trivial on the center, and the exact
sum $|U|^{-1}\sum_x\chi(x^2)$ is one.  This independently
recovers its symmetric invariant form.  The displayed lower
bounds are also checked row by row; in particular no unlisted
degree or quarter-bound exception is omitted.

It remains to pass from a table element to an arbitrary linear
normalizer, without assuming an unproved description of that
normalizer.  The kernel of a nonlinear irreducible
representation of the quasisimple $U$ is central.  An
automorphism of its represented image induces an automorphism
of $S$, and lifts uniquely to the universal perfect cover.
An automorphism trivial on $S$ is trivial on the image: its
central discrepancy is a homomorphism from a perfect group to
an abelian center.  Thus a linear normalizer is, modulo
scalars, inner or induces the one nontrivial outer automorphism.

In the inner case it is a scalar multiple of a represented
element of $U$, covered by the first table.  In the outer case
choose $t\in A\setminus U$ and make an inner adjustment so
that the normalizer $T$ induces precisely conjugation by $t$.
Writing $t^2=u_0\in U$, Schur's lemma gives
$T^2=c\rho(u_0)$.  Rescaling $T$ by a square root of
$c^{-1}$ makes it satisfy the conjugation and power relations
of this fixed $A$.  We obtain an irreducible representation
of $A$ extending $\rho$, hence a row in the checked list.
This argument also covers a nonsplit outer extension and any
represented central quotient.  Scalar rescaling leaves $\nu$
unchanged, proving (SPW.12).

The retained input is \texttt{sp4\_3\_character\_inputs.json},
SHA-256
\texttt{0604b7cbfdd2011c21e46a009674adda4292f1452afbce38ab1593b38c3f388c}.
The complete exporter and checker are
\texttt{export\_sp4\_3.g} and \texttt{check\_sp4\_3.py};
the output \texttt{check\_sp4\_3.json} retains every checked
row and all extremal witnesses.  The false degree-five
quarter-bound expectation is explicitly retained as refuted,
rather than silently discarded.
\end{proof}

\subsection*{D-SPW-04b. The transitive central-product core}

\begin{theorem}[D-SPW-04b: reduction to single degrees four and five]
\label{D-SPW-04b}
Suppose an inducing image $J\subset\operatorname{GL}(V)$
of an actual rank-nine Weil sixfold source has an irreducibly
acting normal central product $N=N_1\cdots N_t$, with all
$N_i/Z(N_i)\simeq S$ and with $J$ transitive on the factors.
Assume its inducing character field contains the original
imaginary quadratic field $K$, as it does for the induced
sources considered here.  Then $t=1$ and
\[
                          \dim V\in\{4,5,6\}.             \tag{SPW.15}
\]
Every tensor case $t\ge2$ is excluded, including degree-six
factors.

If $K$ is one of the eight non-Eisenstein fields in the main
nonsplit problem and the actual original polarization is
nonsplit, the single degree-six case is also excluded by
D-BI-16.  The only remaining possibilities are therefore the
single degree-four and degree-five cases, already subject to
the $G_{32}$ and $G_{33}$ low-degree source restrictions.
No existence or dominance assertion is made for those endpoints.
\end{theorem}

\begin{proof}
As an $N$-module,
$V=V_1\otimes\cdots\otimes V_t$, with irreducible nonlinear
factor modules.  Transitivity of $J$ implies that all factor
degrees have the same value $a\ge4$.  Put $e=a^t$.
The tensor matrix algebras of the intrinsic simple factors
are permuted by $J$.  Hence each normalizing operator is a
tensor product of intertwiners followed by a factor permutation.
For a factor-preserving operator, the largest eigenspace
fraction is at most the largest such fraction of any of its
nonscalar factors.  A nontrivial weighted factor cycle has
largest eigenspace fraction at most $3/4$, by the exact
constant-word argument in D-SPOR-03.  These statements allow
arbitrary scalar phases and compatible outer automorphisms.

The actual intermediate-cover Euler identity and period
factorization give
\[
             h=0,\qquad r\ge12,\qquad D=2e+6,              \tag{SPW.16}
\]
where $r$ counts genuine branch points.  Indeed at genus one
there are at most six genuine branches and at most six
marked-base parameters; at genus at least two the identity
$D=6+2e(1-h)$ is negative since $e\ge4$.  Genuine scalar
inertia has moving rank $e$, so all bounds below apply to it
as well.

If $a=4$ or $a\ge7$, (SPW.12) and the tensor argument give
$\nu(g)\ge e/4$ for every nonscalar inertia.  Equation
(SPW.16) implies $3e\le2e+6$, so $e\le6$.  Thus $a\ge7$
is impossible, and $a=4$ forces $t=1$.
If $a=5$, the corresponding fifth bound gives
$12e/5\le2e+6$, hence $e\le15$.  Therefore $t=1$ here too.

It remains to exclude $a=6$, $t\ge2$; the sixth bound alone
would not do so.  The unique degree-six core character has
indicator $+1$.  Taking the tensor product of its symmetric
forms gives a nondegenerate $N$-invariant symmetric form $B$
on $V$.  Irreducibility makes its invariant line unique, so
normality of $N$ gives a multiplier character
\[
                  B(gv,gw)=\mu(g)B(v,w),\qquad g\in J.
\]
This conclusion does not require a chosen rational model or
an asserted full normalizer identification.

The multiplier is nontrivial.  If it were trivial, every
finite-order element of $J$ would have eigenvalues paired
as $\lambda,\lambda^{-1}$ by $B$; these are complex
conjugates.  Its trace would be real, making the inducing
character field real, contrary to its containing the given
imaginary quadratic $K$.
Because the inertia tuple generates $J$ and has product one,
at least two inertia multipliers are nontrivial.  For each
such inertia $g$, its fixed space is totally isotropic for
$B$, since
$(\mu(g)-1)B(v,w)=0$ on that space.  Its moving rank is
therefore at least $e/2$.
Every other inertia has moving rank at least $e/6$ by
(SPW.12) and the tensor-cycle bound.  Consequently
\[
             2e+6=D\ge 2(e/2)+10(e/6)=8e/3,
             \qquad e\le9.                               \tag{SPW.17}
\]
This contradicts $e=6^t\ge36$ when $t\ge2$.
Only the single degree-six representation remains, proving
(SPW.15).

Under the final nonsplit, non-Eisenstein hypotheses, D-BI-16
applies to any degree-six inducing representation, independently
of its inducing field.  Its sole allowed faithful image is
$\mu_4G_{34}$.  This has projective order
$39191040/6=6531840$, whereas the present normalizer has
projective order at most $|\operatorname{Aut}(S)|=51840$.
Thus the degree-six case here cannot be that allowed image.
Equivalently, it is the already excluded real $E_6$ case.
The degree-four and degree-five conclusions refer back to
their previously proved low-degree restrictions and assert no
additional normalizer identification.  This completes the
stated reduction.
\end{proof}
